\label{sec:conclusion}

We built a Mininet-based harness to compare single-path TCP with a multipath-capable Linux configuration under emulated mobility stressors. Independent TCP baselines track configured link capacities well (Experiment~1). Running two interface-bound \texttt{iperf3} flows shows that both emulated paths can be utilized together, with combined throughput matching or slightly exceeding the sum of sequential baselines in our logs (Experiment~2). A hard interface failure on one path stops that subflow immediately, but the surviving path continues delivering throughput without restarting the server (Experiment~3), which is the operational notion of continuity we could measure. Bandwidth collapse on one path reduces aggregate throughput but the stable path absorbs additional load, recovering a substantial fraction of the pre-event aggregate (62.5\% retained in Experiment~4). An RTT spike on a single path devastates that TCP flow (13.5\% retained on Path~2 alone), whereas the dual-path setup retains 71.0\% of its pre-spike combined average (Experiment~5). Experiment~6 steps up delay on the initially preferred path: aggregate throughput before and after the handover window differs only modestly ($24.89$ vs.\ $23.39$~Mb/s in our log), with Path~1 incurring most retransmissions---consistent with gradual load shift under scheduler control rather than catastrophic outage.

These results support the intuition behind MPTCP in heterogeneous access: spare paths provide redundancy against localized impairment. Our study does not yet quantify Jain fairness across competing flows; parallel \texttt{iperf3} clients remain an approximation of full multipath socket behavior in every scenario---so claims should be read as emulator-scoped evidence rather than universal WAN performance.
